\chapter*{سپاسگزاری}
\addcontentsline{toc}{chapter}{سپاسگزاری}
با درود و ستایش بی‌پایان، سخن را با نام آن یگانۀ بی‌مانند می‌گشایم که الطاف بی‌حدوحصرش رهگشای گام‌های ناتوان من در این وادی دانش بود و مرا به فرجام این مقصود موفق و مؤید گردانید.

پس از ستایش ایزدی، ارج نهادن به مقام والای استاد فرزانه و راهنمای خویش، سرکار خانم دکتر زیبا اسلامی، برترین فریضۀ سپاسگزاری است. آن بزرگوار با شکیبایی کم‌مانند، ژرفای بینش و رهنمودهای سودمندشان همواره روشنی‌بخش مسیر این پژوهش بوده‌اند و بی‌یاری خردورزانۀ ایشان، نگارش این اوراق میسر نمی‌شد. امید که این قلیل بضاعت پاسخی هرچند ناچیز به رنج‌های پیگیر آن استاد گران‌مایه باشد.

در ادامه، مراتب امتنان و سپاس صمیمانۀ خویش را به داوران گرامی این پایان‌نامه، جناب آقای دکتر هادی فراهانی و جناب آقای دکتر سعید رضا خردپیشه، تقدیم می‌دارم. دقت نظر، ارزیابی علمی و پیشنهادهای ارزشمند ایشان در فرایند داوری، در اصلاح روش پژوهش، روشن‌تر شدن گزارش نتایج و ارتقای کیفیت علمی و نگارشی این پایان‌نامه نقشی بسزا داشته است. از وقتی که با بزرگواری به مطالعه و بررسی این پژوهش اختصاص دادند، صادقانه سپاسگزارم.

بی‌تردید، پیمودن بخش‌های عملی این تحقیق جز با یاری دست‌اندرکاران عرصۀ سخت‌افزار میسر نبود. از جناب مهندس ناظری، از دانشگاه فردوسی مشهد، که با گشاده‌دستی تجهیزات و زیرساخت‌های سخت‌افزاری لازم را در دسترس اینجانب نهادند، صمیمانه سپاسگزارم. همچنین از جناب مهندس شمس‌آبادی و همسر ایشان، مهندس نازفروش، که در روزگار دانشجویی همواره در گشایش دشواری‌ها و ارائۀ یاری‌های کارگشا یاور من بودند، از صمیم قلب قدردانی می‌کنم. محبت‌های بی‌چشمداشت این بزرگواران هیچ‌گاه از خاطر نخواهد رفت.

در پایان، بر خود می‌بالم که در این مسیر پرفرازونشیب با بزرگوارانی همسفر بودم که یادشان سرمایۀ دل است. از درگاه خداوند متعال، دوام تندرستی و فزونی توفیق همگان را مسئلت می‌کنم.

\begin{flushleft}
سجاد رنجبر یزدی\\
شهریور ۱۴۰۵
\end{flushleft}
